\section{Conclusions and Outlook}
\label{sec:conclusions}

Version 4.2 adopts the factor-12 upper length-scale setting for the full 2016
sample and completes the corresponding three-campaign result. The search
intervals and wider GP supports remain 19--90 and 14--135~MeV for 2015, 39--180
and 30--210~MeV for 2016, and 50--250 and 40--300~MeV for 2021. The full 2015
and 2016 samples and the 2021 10\% sample enter a shared-$\eps^2$ likelihood.

The adopted setting follows the controlled factor-8, 10, 12, 15, and 20 range
study. The v4 ceiling is occupied at all 142 full-2016 hypotheses; factor 10
remains occupied at 56, while factors 12, 15, and 20 have no occupied
hypotheses. Factor 12 is the first nonbinding value and is followed by a stable
plateau in log marginal likelihood, observed yield limits, and local $Z$.
Repeated actual fits resolve the reviewed branches, and no GP state, limit, or
$p$-value is interpolated. The rule does not select whichever observed limit is
tighter or whichever local $p_0$ is smaller.

The exact 415-state reconstruction is finite on all 232 combined hypotheses.
For the combined search, 300 shared mass-local conditional background-only
pseudoexperiments provide the median and central limit quantiles. They retain the
reviewed GP states without refitting and use the same asymptotic inner 90\%
\CLs{} construction as the observed limit. The observed-to-median ratio ranges
from 0.514396 at 73~MeV to 3.29400 at 65~MeV. The corresponding mass-local tail
counts reach $0/300$ on the strong side at 71--73~MeV and on the weak side at
21, 65, and 66~MeV; a zero count is finite ensemble resolution, not an exact
probability.

The individual full-2015, full-2016, and 2021 10\% observed limits are reported
alongside the combined result. Auxiliary 100-pseudoexperiment conditional bands
document each standalone likelihood and each pairwise shared-$\eps^2$
combination. They select the first 100 entries from the accepted
300-pseudoexperiment parent stream, reuse campaign draws across scopes, and are
not used as direct-coverage or scan-wide discovery ensembles. The smallest
campaign-level asymptotic $p_0$ values are
$8.46283\times10^{-4}$ at 51~MeV, $3.08781\times10^{-4}$ at 90~MeV, and
$1.05702\times10^{-5}$ at 65~MeV, respectively. The combined minimum is
$3.25918\times10^{-5}$ at 65~MeV ($Z=3.99321$). With
$N_{\mathrm{eff}}=35.3814$, the fixed-card analytic \v{S}id\'ak reference is
$p=0.00115250$ ($Z=3.04783$).

The exact 65~MeV extraction gives
$\widehat{\eps^2}=(6.42706\pm1.61009)\times10^{-6}$. Its dataset-block
decomposition shows that the localized feature is predominantly carried by the
2021 10\% sample, is modestly supported by 2015, and is opposed by 2016 at the
common-coupling yield. The displayed background-subtracted spectra and profiled
residuals document this structure without treating individual bins as independent
significances.

The August 6 follow-ups leave that accepted result unchanged.  Exact common
0.625~MeV inference bins cannot be constructed from the histogram-only 2015 and
2016 inputs; the finer nearest source-compatible choice that retains all three
full supports is 0.5~MeV, and its refit gives
$Z_{\mathrm{local}}=3.85429$, a fixed-mass shift of $-0.13892$, with an
essentially unchanged Wald scale.  The repaired coefficient display identifies
the former negative 2016 endpoint as a signed-Wald convention and instead shows
the nominal physical profile set $[0,4.08001]\times10^{-6}$.  The
observed-equivalent 2021-statistics response can be numerically lower than the
published BaBar visible contour, but it preserves current fluctuations and is
not an expected sensitivity or a future full-data result.  Finally, both
fixed-seed background replacements remove the positive estimate specifically
at 65~MeV; their full-scan boundary excursions remain conditional local
diagnostics rather than global-null claims.  The lower-mass shoulder in the
functional-form draw is not purely statistical: the deterministic functional
mean gives $Z=2.268$ at 61~MeV under the GP analysis, while the particular
draw moves the maximum to 62~MeV and raises it to $Z=3.690$.  In the separate
fixed-GP-mean ten-draw screen, the 2.25- and 2.5-$\sigma_m$ requests are the
same sixteen production bins.  The 3-$\sigma_m$ replacement gives a paired
median 65~MeV limit ratio of 1.17957 relative to that common narrow scope; the
reported ribbons are descriptive conditional spreads, not expected bands.

The conditional limit ensemble does not establish direct \CLs{} coverage, refit
the GP, or supply a scan-wise discovery calibration. The analytic \v{S}id\'ak
reference likewise does not account for the observed-data upper-bound study.
Factor 12 is accepted for the v4.2 result, while the post-selection and coverage
qualifications remain explicit.

The v4.1 source/exposure length-scale pilot and its v4.5 ten-toy refmatched
continuation are now preserved in Appendix~\ref{app:v4p5-ten-toy-source-exposure}.
Version 4.6 analyzes the complete 100-toy $1\%\times10$, native-10\%,
$1\%\times100$, and $10\%\times10$ ensembles under the frozen v4.2 2021 card.
All 100 background pseudo-histograms are rerun uniformly at 65, 90, 120, 180,
and 210~MeV with 0-, 1-, 3-, and 5-$\sigma_A$ injections; old extraction rows
are not pooled with the new estimator output.

The predeclared numerical gate separates the pseudoexperiment stream from the
optimizer-start stream and gives each GP fit a nominal start plus 12 randomized
restarts.  Reference fits must reproduce their highest-LML branch, while nonzero
fits receive adaptive independent attempts only after pull- and amplitude-blind
technical triggers.  Selection never uses fitted amplitude, pull sign or size,
$\eps^2$, or proximity to the injected value.  Six irreproducible state keys are
excluded from 8000, but no entire background-toy index is discarded and every closure
cell retains 99 or 100 accepted toys.  Statistical extremes that pass the pull-blind
technical gate remain in the arithmetic moments; repeat-triggered factor-15 boundary
solutions are retained when replicated.
The frozen covariance repair is also audited explicitly: its raw-matrix tolerance is
a repairability criterion, not a positive-semidefiniteness claim, and a minimal
diagonal-loading alternative changes any reconstructed pull by less than 0.003.

This policy distinguishes two effects that the ten-toy screen could not.  The
dramatic $10\%\times10$ branches at 65 and 180~MeV were numerical: repeated
maximum-LML fits recover stable ordinary-uncertainty solutions without removing
their background toys.  In contrast, the pilot-selected $1\%\times10$, 65~MeV
zero-signal offset persists in the independent toys-10--99 reserve, with
$\bar p=0.768$, width 0.944, and 95\% interval $[0.570,0.965]$.  The full-100
mean is 0.789, and the deterministic analytic stress-truth mean gives 0.837.
This is a local functional-form/GP closure offset, not a few bad toys.  Native
10\% at 65~MeV and $10\%\times10$ at 180~MeV also retain exploratory
Holm-corrected material zero-signal offsets.  All are conditional on the smooth
stress truth and do not measure a bias of the observed spectrum.

Pull widths otherwise lie between 0.879 and 1.237 across accepted cells.  The
single value above 1.2 is the native-10\%, 65~MeV, $z=5$ width of 1.237.
Factor-15 upper-bound contacts remain absent in the lower-exposure ensembles but
occur in 85/1999 accepted $1\%\times100$ fits and 134/1995 accepted
$10\%\times10$ fits.  They are reported rather than hidden, and the full-100
study does not retroactively alter the frozen candidate card.

The source functions are independently fitted and the cross-source ensembles are
unpaired.  The expected normalization differs by a factor 1.1296466 in both
requested comparisons, so the signed $\eps^2$ distributions combine
source/selection shape and normalization rather than a pure exposure effect.
They remain extraction-coordinate diagnostics, not expected-limit bands.
Full-100 conditional injection--extraction closure is now documented, but direct
\CLs{} coverage, independent alternate-truth closure, and a scan-wise
background-only maximum-$q_0$ ensemble remain separate requirements.

Version 4.7 tests whether a source-dedicated description of the 2021 turn-on
removes the near-threshold closure response.  Independent fits to the native
1\% and 10\% spectra over 35--80~MeV select a logistic turn-on times an
exponentiated degree-five Chebyshev series as the lowest common candidate
passing fixed native-bin and five-bin goodness-of-fit gates.  Twenty Poisson
background spectra per source-derived scenario are tested at 55, 60, 65, and
70~MeV with 0-, 1-, 3-, and 5-$\sigma_A$ injections.  The v4.2/v4.5
40--300~MeV support, log preprocessing, factor-1.1/factor-15 2021 kernel
bounds, and $2.25\sigma_m$ blind/training geometry remain fixed.

The dedicated truth gives partial mitigation rather than uniform closure.  At
65~MeV the $1\%\times10$ zero-signal mean changes from the unpaired archived
v4.6 value 0.789 to $0.300\;[-0.219,0.819]$, and the corresponding analytic
mean response changes from 0.837 to 0.252.  Native 10\% remains at
$0.641\;[0.082,1.199]$, compared with the unpaired v4.6 mean 0.716, and its
analytic response changes from 0.723 to 0.563.  Both scenarios have negative
55~MeV means, $-0.985$ and $-0.824$, while $1\%\times10$ has a positive
70~MeV mean of 0.625.  Pull widths across all 32 cells range from 0.644 to
1.196.  Each interval contains only 20 backgrounds and cells share their
background streams; these values are conditional diagnostics, not coverage or
independent repeated confirmations.

The initial complete v4.7 optimizer ledger exposed three additional displaced
single-attempt branches through reference-relative GP likelihood and kernel
coordinates.  A pull-blind repeat gate was amended after that audit and all 40
tasks were uniformly rerun as v4.7.1.  The final ledger contains all 640 unique
states with zero exclusions, 14 repeat-triggered injected fits, no kernel-bound
contacts, strictly positive training counts, and a smallest raw covariance
eigenvalue ratio of $-1.09\times10^{-6}$.  The amendment is explicitly
post-audit rather than prospectively predeclared.  Superseded rows are retained
only as numerical provenance and are not pooled with the accepted summaries.

The truth construction also has declared limits.  Its join rescales the full
tail above 80~MeV by about 1.3\%, and the native-10\% 80--300~MeV tail has
deviance/ndf 2.861.  The full-range Gaussian injection leaves a median 2.53\%
of realized signal in the retained training bins.  The archived old 100-toy
65~MeV ensemble is unpaired with the new 20-toy streams, so the comparison does
not isolate a causal effect of the functional-form replacement.  Follow-ups
should consequently freeze and compare a raised-support crop of the cached
toys, a lower-support study with newly generated and validated sub-40-MeV bins
but identical bins above 40~MeV, discretely effective blind/training-window
variations, and, if still required, a $k_{\min,2021}=1.0$ test on matched
pseudo-data.  A parametric turn-on mean plus GP residual or a
nonstationary/input-warped covariance should be qualified with held-out fake
gaps and a predeclared spurious-signal requirement before use.  A validated
set of surviving analytic families could instead enter as a discrete-profiled
function-choice nuisance.  None of these recommendations changes the accepted
v4.2 card or observed result.

Version 4.8 carries out that residual-structure follow-up with two deliberately
restricted source models.  Blocked source-only folds select the fixed K2
natural log spline over the K3 comparator, while the second model uses fixed
quintic overlaps between low-, middle-, and high-mass components.  Neither
model passes the frozen source-qualification gates.  The K2 native-1\% and
native-10\% conservative tangent projections reach 0.245 and 0.237, and the
corresponding regional values reach 0.631 and 0.328, all above the 0.04 gate.
The fake-gap and deterministic injected-refit gates also fail.  A small number
of knots or regions therefore does not establish signal rigidity.

The requested conditional extraction screen is nevertheless complete.  Each
model has 20 backgrounds in native 1\%, $1\%\times10$, $1\%\times100$,
native 10\%, and $10\%\times10$ at five masses and four injection strengths:
2000 raw and accepted states per model, with no exclusions.  The exploratory
zero-signal screen flags three of 25 K2 cells, all at 65~MeV, and sixteen of 25
regional cells.  The regional sign-changing offsets align with its two overlap
zones and reach 7.558 mean-pull units.  After subtracting the matched
zero-injection estimate, the all-row paired-response medians are 0.9804 for K2
and 0.9708 for the regional blend, but rare accepted branch changes give much
larger losses.  These downstream responses occur only after the generator has
been frozen and cannot override the adverse source-fit influence audit.

The original pull-blind pilot's factor-25 fallback is insufficient for the K2
native-1\% lane: all nine pilot states and 345 of 400 closure states contact
that ceiling.  A post-closure addendum uses eight fresh background-only toys
and no inference quantities.  Factor 35 fails at 120~MeV; factor 50 passes
against a factor-75 sentinel at 65, 120, and 210~MeV in both the three-toy
selection and five-toy confirmation sets.  This qualifies only the targeted
background-optimizer ceiling diagnostic.  It neither reruns nor uncensors the
factor-25 closure, establishes a common all-lane ceiling, nor changes the
production v4.2 factor-15 card.

Accordingly, K2 is retained only as the less pathological requested
conditional stress truth, and the regional blend is not suitable for
quantitative closure claims.  Neither is a qualified generator, a coverage
result, an expected band, observed-data bias, a limit, or an exclusion.  The
accepted v4.2 search result is unchanged.

Production-faithful functional-form closure and direct limit coverage should be
rerun under a card frozen independently of observed limit direction or local
significance. A separate scan-wise
background-only maximum-$q_0$ ensemble would be required for a toy-calibrated global
discovery statement. The eventual full-2021 observed spectrum remains a new analysis
input, not a deterministic scale factor applied to the present 10\% observation.
